\section{Collapsing a probability component}

The measure-theoretic construction in Example~3.7.2(3) admits the following
general form. It will be applied with the finite places as $A$ and the
transcendental points of the unit circle as $B$.

\begin{lemma}[Integration on a disjoint union]
\label{lem:sum-measure-integration}
\leanok
\source{chen-moriwaki-adelic:page-0220}
\lean{ArakelovGeometry.ProbabilityCollapse.integrable_sumMeasure_iff,ArakelovGeometry.ProbabilityCollapse.ae_sumMeasure_iff,ArakelovGeometry.ProbabilityCollapse.integral_sumMeasure}
Let $(A,\mu)$ and $(B,\nu)$ be measure spaces, and equip $A\amalg B$ with
$\lambda=(\iota_A)_*\mu+(\iota_B)_*\nu$.
A real function $f$ is integrable for $\lambda$ if and only if both
$f_A=f\circ\iota_A$ and $f_B=f\circ\iota_B$ are integrable, and then
$\int f\,d\lambda=\int f_A\,d\mu+\int f_B\,d\nu$.
An equality of functions holds $\lambda$-almost everywhere if and only
if it holds almost everywhere on both components.
\end{lemma}
\begin{proof}
Each inclusion is a measurable embedding. Pullback along it identifies
measurable representatives and null sets for its image measure with those
on the corresponding component. A set has zero measure for a sum of two
positive measures precisely when it has zero measure for both. Similarly,
the nonnegative integral of $|f|$ for the sum is the sum of the two
nonnegative integrals. This is finite precisely when both summands are
finite, proving the integrability assertion. Applying additivity to the
positive and negative parts gives the claimed real integral formula.
\end{proof}

\begin{lemma}[Collapsing a probability space]
\label{lem:probability-collapse}
\leanok
\source{chen-moriwaki-adelic:page-0220}
\uses{lem:sum-measure-integration}
\lean{ArakelovGeometry.ProbabilityCollapse.measurePreserving_collapse,ArakelovGeometry.ProbabilityCollapse.averageL1LinearMap,ArakelovGeometry.ProbabilityCollapse.averageL1_toL1,ArakelovGeometry.ProbabilityCollapse.integral_averageL1,ArakelovGeometry.ProbabilityCollapse.averageL1_compMeasurePreserving}
Assume $\nu(B)=1$. Give $A\amalg\{*\}$ the measure
$\lambda_0=(\iota_A)_*\mu+\delta_*$.
The map $q:A\amalg B\to A\amalg\{*\}$ which fixes $A$ and sends
$B$ to $*$ is measurable and satisfies $q_*\lambda=\lambda_0$.
The formula
\[
 Jf(a)=f(a)\quad(a\in A),\qquad
 Jf(*)=\int_B f_B\,d\nu
\]
induces a real-linear map $J:L^1(\lambda)\to L^1(\lambda_0)$.
It preserves the integral and satisfies $Jq^*=\operatorname{id}$.
\end{lemma}
\begin{proof}
The restrictions of $q$ to the two components are respectively the
inclusion of $A$ and a constant map, so $q$ is measurable. Its pushforward
is $(\iota_A)_*\mu+\nu(B)\delta_*=\lambda_0$.
For an integrable $f$, the preceding lemma shows that $f_A$ is integrable
and that the integral defining $Jf(*)$ exists. Every real constant on the
one-point probability space is integrable. Thus $Jf$ is integrable.
If $f=g$ almost everywhere, the componentwise criterion gives
$f_A=g_A$ almost everywhere and $f_B=g_B$ almost everywhere. The latter
equality implies equality of their integrals, proving that $J$ is
well-defined on classes. Additivity of the integral on integrable
functions and homogeneity prove real linearity. The disjoint-union formula
gives
\[
 \int Jf\,d\lambda_0
 =\int_A f_A\,d\mu+\int_B f_B\,d\nu
 =\int f\,d\lambda.
\]
For a function $g$ on $A\amalg\{*\}$, its pullback restricts to $g$
on $A$ and the constant $g(*)$ on $B$. Since $\nu(B)=1$, averaging this
constant gives $g(*)$. Therefore $J(g\circ q)=g$ pointwise, which gives
the asserted identity on integrable equivalence classes.
\end{proof}
